% chapters/ch02.tex
\chapter{Marco teórico}

\begin{refsection}

\section{Inclusión educativa en la sociedad del conocimiento}

\subsection{Evolución del concepto de inclusión educativa}

La inclusión educativa ha evolucionado desde enfoques centrados en la integración de estudiantes con necesidades educativas especiales hacia una concepción más amplia que reconoce la diversidad como una condición inherente de los sistemas educativos. En esta perspectiva, la inclusión no se reduce a la presencia física en el aula, sino que implica participación activa, aprendizaje significativo, reconocimiento de diferencias individuales y transformación de las culturas, políticas y prácticas institucionales \parencite{ainscow2006,unesco2020}.

En la sociedad del conocimiento, los procesos educativos se encuentran atravesados por tecnologías digitales, flujos masivos de información y nuevas formas de producción y circulación del conocimiento \parencite{castells2010}. En este escenario, la inclusión demanda estrategias sistemáticas de diagnóstico y seguimiento que permitan identificar, con oportunidad y precisión, barreras para el aprendizaje y la participación, así como fortalezas y apoyos requeridos, evitando enfoques reduccionistas basados únicamente en el rendimiento.

\subsection{Inclusión educativa y equidad}

La inclusión se articula estrechamente con la equidad, entendida como la provisión de apoyos diferenciados para garantizar oportunidades reales de aprendizaje. En consecuencia, la equidad no equivale a tratar a todos por igual, sino a reconocer desigualdades estructurales y responder con intervenciones pedagógicas y organizacionales basadas en evidencia \parencite{oecd2018}.

Un reto central consiste en identificar necesidades y apoyos sin derivar en prácticas estigmatizantes. La literatura sobre pedagogía inclusiva enfatiza que la diversidad debe abordarse desde el diseño de experiencias de aprendizaje flexibles y desde la eliminación de barreras, más que desde la categorización rígida del estudiantado \parencite{florian2011}. En este punto, la tecnología educativa basada en datos abre oportunidades relevantes, siempre que se fundamente en principios pedagógicos, de accesibilidad y de ética de datos.

\section{Diversidad, discapacidad y barreras para el aprendizaje}

\subsection{Enfoques contemporáneos sobre discapacidad y aprendizaje}

En investigación educativa inclusiva resulta clave diferenciar entre (i) condiciones o características personales y (ii) barreras del entorno que limitan el acceso, la participación y el aprendizaje. En enfoques contemporáneos, la discapacidad se comprende como una interacción entre condiciones de salud y factores contextuales (ambientales y personales), lo cual desplaza la mirada exclusiva del déficit y orienta la atención hacia ajustes razonables, apoyos y accesibilidad \parencite{whoicf2001}.

Desde esta perspectiva, la diversidad del estudiantado puede involucrar, entre otras, variaciones sensoriales (p.\ ej., visuales o auditivas), motrices, comunicativas, cognitivas, socioemocionales y del neurodesarrollo, así como condiciones derivadas de contextos socioculturales y lingüísticos. Para fines pedagógicos y diagnósticos inclusivos, el énfasis no se sitúa en etiquetar, sino en reconocer necesidades de apoyo, preferencias de aprendizaje, barreras de acceso y condiciones contextuales que inciden en la trayectoria educativa.

\subsection{Barreras para el aprendizaje y la participación}

Las barreras pueden manifestarse en múltiples niveles: curriculares (rigidez de contenidos y evaluación), didácticos (escasa variabilidad de estrategias), comunicativos (formatos no accesibles), tecnológicos (plataformas no navegables con teclado o lector de pantalla), organizacionales (falta de rutas de apoyo) y actitudinales (bajas expectativas, estigmas). En consecuencia, un diagnóstico inclusivo debe integrar información pedagógica y contextual, además de evidencias del desempeño, para orientar decisiones basadas en apoyos y oportunidades.

\section{Accesibilidad digital y diseño inclusivo en plataformas educativas}

\subsection{Accesibilidad como condición de inclusión mediada por tecnología}

En la sociedad del conocimiento, una proporción significativa de las oportunidades educativas se materializa mediante entornos digitales (LMS, plataformas, recursos interactivos, repositorios y videoconferencias). Por tanto, la inclusión requiere que estos entornos sean accesibles, usables y comprensibles para personas con diversos perfiles funcionales y contextos de uso.

En el ámbito web, las Pautas de Accesibilidad para el Contenido Web (WCAG) establecen criterios ampliamente aceptados para garantizar que los contenidos sean \textit{perceptibles, operables, comprensibles y robustos} \parencite{w3cwcag}. Complementariamente, la accesibilidad puede alinearse con estándares y marcos técnicos que orientan requisitos para productos y servicios TIC \parencite{etsi301549}, así como con especificaciones de accesibilidad semántica para componentes e interfaces \parencite{w3caria}.

\subsection{Diseño centrado en el usuario y accesibilidad desde el inicio}

Desde la ingeniería de sistemas y la interacción humano--computador, el diseño centrado en el usuario (o diseño centrado en las personas) plantea que la calidad del sistema depende de comprender a los usuarios, sus tareas, contextos y limitaciones, incorporando iteración y evaluación sistemática a lo largo del ciclo de vida \parencite{iso9241210}. En investigación educativa aplicada, esto resulta especialmente pertinente cuando la plataforma se orienta a poblaciones diversas y a contextos institucionales con brechas de conectividad, equipamiento y alfabetización digital.

En coherencia con este enfoque, una plataforma orientada al diagnóstico inclusivo debe integrar, al menos, los siguientes principios operativos:

\begin{itemize}
\item \textbf{Acceso multimodal:} alternativas textuales para contenido no textual, subtitulado y transcripción, control de velocidad y formato.
\item \textbf{Operabilidad:} navegación por teclado, foco visible, estructura consistente y prevención de bloqueos por tiempo.
\item \textbf{Comprensibilidad:} lenguaje claro, instrucciones explícitas, retroalimentación oportuna y reducción de carga cognitiva.
\item \textbf{Robustez:} compatibilidad con tecnologías de asistencia (lectores de pantalla), marcado semántico adecuado y componentes accesibles.
\end{itemize}

La accesibilidad no es un módulo aislado: constituye una condición transversal del diseño tecnológico, del modelo de datos (por ejemplo, el registro de preferencias de acceso) y de la visualización de resultados diagnósticos (por ejemplo, reportes legibles, exportables y configurables).

\section{Diseño Universal para el Aprendizaje (DUA)}

\subsection{Fundamentos del DUA}

El \DUA\ surge como un marco pedagógico orientado a anticipar la diversidad del alumnado, en lugar de reaccionar a ella mediante adaptaciones posteriores \parencite{cast2018,meyer2014}. Este enfoque se apoya en hallazgos sobre redes implicadas en el aprendizaje y propone organizar el diseño pedagógico desde tres principios:

\begin{itemize}
\item proporcionar múltiples medios de compromiso;
\item proporcionar múltiples medios de representación;
\item proporcionar múltiples medios de acción y expresión.
\end{itemize}

El \DUA\ no se limita a incorporar recursos variados, sino que promueve un diseño intencional que reduzca barreras, diversifique rutas de acceso al aprendizaje y favorezca la participación, la autorregulación y el progreso.

\subsection{DUA y diagnóstico educativo}

Si bien el \DUA\ ha sido ampliamente abordado en diseño curricular y práctica docente, su implementación efectiva requiere diagnósticos que permitan conocer, de manera sistemática, necesidades de apoyo, preferencias, barreras y condiciones de acceso. En contraste, los diagnósticos tradicionales suelen basarse en pruebas estandarizadas, centrarse en déficits y ofrecer información limitada para decisiones pedagógicas contextualizadas \parencite{meyer2014}.

La integración de tecnologías digitales permite avanzar hacia diagnósticos dinámicos y continuos, mediante la triangulación de cuestionarios, observación, evidencias de desempeño y trazas digitales. De esta forma, el diagnóstico puede alinearse con el \DUA\ al generar información accionable sobre apoyos y estrategias, evitando enfoques reduccionistas.

\section{Educación basada en datos y analítica del aprendizaje}

\subsection{Analítica del aprendizaje y toma de decisiones pedagógicas}

La educación basada en datos se refiere al uso sistemático de información educativa para orientar decisiones, mejorar procesos y evaluar intervenciones \parencite{mandinach2016}. En este marco, la analítica del aprendizaje se entiende como la medición, recopilación, análisis y reporte de datos sobre estudiantes y contextos con el propósito de comprender y optimizar el aprendizaje \parencite{siemens2012}.

Sin embargo, el valor pedagógico de la analítica depende de su interpretabilidad y de su articulación con el conocimiento docente. La evidencia educativa debe convertirse en conocimiento útil para decisiones de aula, no en métricas descontextualizadas. Por ello, los sistemas inteligentes orientados al diagnóstico inclusivo requieren: (i) modelos de indicadores pedagógicamente significativos, (ii) explicaciones comprensibles para docentes y directivos y (iii) mecanismos de retroalimentación para iterar el modelo.

\subsection{Riesgos y desafíos éticos de la educación basada en datos}

El uso de datos educativos plantea riesgos asociados a privacidad, vigilancia, sesgos algorítmicos y exclusión digital. En particular, cuando los modelos se aplican a poblaciones diversas, existe el riesgo de amplificar desigualdades si se usan variables \textit{proxy} de vulnerabilidad sin salvaguardas o si se interpretan correlaciones como causalidades.

La \textcite{unesco2021} enfatiza que los sistemas basados en \IA\ deben diseñarse bajo principios de justicia, transparencia, explicabilidad, responsabilidad y no discriminación, especialmente en ámbitos sensibles como la educación. En consecuencia, esta investigación asume que la educación basada en datos no es neutral y que su implementación debe orientarse a fortalecer la inclusión, evitando reproducir desigualdades preexistentes.

\section{Articulación conceptual: inclusión, DUA, accesibilidad y educación basada en datos}

La convergencia entre inclusión educativa, \DUA, accesibilidad digital y educación basada en datos constituye el núcleo teórico de esta investigación doctoral:

\begin{itemize}
\item la inclusión educativa define el \textbf{fin}: equidad, participación y eliminación de barreras;
\item el \DUA\ orienta el \textbf{cómo pedagógico}: diseño anticipatorio para la diversidad;
\item la accesibilidad digital establece la \textbf{condición de posibilidad} en entornos tecnológicos;
\item la educación basada en datos aporta el \textbf{cómo metodológico-tecnológico}: evidencia para comprender y mejorar.
\end{itemize}

Con el fin de sintetizar esta articulación, se presenta en primer lugar una figura de \textbf{arquitectura conceptual}, orientada a mostrar los pilares teóricos que fundamentan la propuesta, y posteriormente una figura de \textbf{arquitectura funcional}, dirigida a explicar la lógica operativa del sistema tecnológico de diagnóstico educativo inclusivo. Esta división permite representar con mayor claridad tanto el sustento teórico del modelo como su traducción en componentes funcionales dentro del sistema.

\begin{figure}[p]
\centering
\begin{adjustbox}{max totalsize={\textwidth}{0.90\textheight},center}
\begin{tikzpicture}[
  font=\small,
  node distance=12mm and 10mm,
  pillar/.style={
    draw,
    rounded corners=2mm,
    align=center,
    inner sep=7pt,
    text width=3.7cm
  },
  core/.style={
    draw,
    rounded corners=2mm,
    align=center,
    inner sep=9pt,
    text width=11.5cm
  },
  block/.style={
    draw,
    rounded corners=2mm,
    align=center,
    inner sep=8pt,
    text width=5.0cm
  },
  wideblock/.style={
    draw,
    rounded corners=2mm,
    align=center,
    inner sep=8pt,
    text width=11.0cm
  },
  arrow/.style={
    -{Stealth[length=2.6mm]},
    line width=0.85pt
  }
]

% ======================================================
% PILARES TEÓRICOS
% ======================================================
\node[pillar] (incl) {\textbf{Inclusión educativa}\\
Equidad\\
Participación\\
Eliminación de barreras};

\node[pillar, right=of incl] (dua) {\textbf{Diseño Universal para el Aprendizaje}\\
Múltiples medios de compromiso\\
Múltiples medios de representación\\
Múltiples medios de acción y expresión};

\node[pillar, right=of dua] (acc) {\textbf{Accesibilidad digital}\\
Perceptibilidad\\
Operabilidad\\
Comprensibilidad y robustez};

\node[pillar, right=of acc] (data) {\textbf{Educación basada en datos}\\
Analítica del aprendizaje\\
Interpretación pedagógica\\
Mejora basada en evidencia};

% ======================================================
% ARTICULACIÓN CENTRAL
% ======================================================
\node[core, below=18mm of dua] (art) {\textbf{Articulación conceptual integradora}\\
Convergencia entre inclusión educativa, \DUA, accesibilidad digital y educación basada en datos\\
\emph{Marco teórico-metodológico para comprender la diversidad del aprendizaje y orientar decisiones pedagógicas basadas en evidencia}};

\draw[arrow] (incl.south) -- ++(0,-5mm) -| ([xshift=-35mm]art.north);
\draw[arrow] (dua.south)  -- ++(0,-5mm) -- ([xshift=-10mm]art.north);
\draw[arrow] (acc.south)  -- ++(0,-5mm) -- ([xshift=10mm]art.north);
\draw[arrow] (data.south) -- ++(0,-5mm) -| ([xshift=35mm]art.north);

% ======================================================
% SISTEMA COMO MATERIALIZACIÓN
% ======================================================
\node[wideblock, below=14mm of art] (sys) {\textbf{Sistema tecnológico inteligente de diagnóstico educativo inclusivo}\\
Propuesta de mediación pedagógica y tecnológica orientada al diagnóstico dinámico, contextualizado, interpretable y éticamente responsable};

\draw[arrow] (art) -- (sys);

% ======================================================
% DIMENSIONES DE MATERIALIZACIÓN
% ======================================================
\node[block, below left=14mm and 14mm of sys] (pedagogic) {\textbf{Dimensión pedagógica}\\
Diagnóstico orientado a apoyos\\
Interpretación docente\\
Recomendaciones alineadas con DUA};

\node[block, below=14mm of sys] (analytic) {\textbf{Dimensión analítica}\\
Integración de datos\\
Construcción de indicadores\\
Identificación de patrones};

\node[block, below right=14mm and 14mm of sys] (ethical) {\textbf{Dimensión ética y de acceso}\\
Privacidad y protección de datos\\
No discriminación\\
Accesibilidad y usabilidad};

\draw[arrow] (sys.south) -- ++(0,-6mm) -| (analytic.north);
\draw[arrow] (sys.south) -- ++(0,-6mm) -| (pedagogic.north);
\draw[arrow] (sys.south) -- ++(0,-6mm) -| (ethical.north);

% ======================================================
% RESULTADO CONCEPTUAL
% ======================================================
\node[core, below=16mm of analytic] (result) {\textbf{Resultado teórico del modelo}\\
El diagnóstico educativo inclusivo se concibe como un proceso continuo, multimodal, contextualizado y basado en evidencia, orientado a reconocer barreras, fortalezas, necesidades de apoyo y oportunidades de mejora pedagógica};

\draw[arrow] (pedagogic.south) -- ++(0,-5mm) -| ([xshift=-30mm]result.north);
\draw[arrow] (analytic.south)  -- (result.north);
\draw[arrow] (ethical.south)   -- ++(0,-5mm) -| ([xshift=30mm]result.north);

\end{tikzpicture}
\end{adjustbox}
\caption{Arquitectura conceptual del diagnóstico educativo inclusivo basado en datos.}
\label{fig:arquitectura-conceptual-marco-teorico}
\end{figure}

La Figura \ref{fig:arquitectura-conceptual-marco-teorico} muestra que el modelo conceptual no surge de un único enfoque, sino de la integración de cuatro pilares complementarios. La inclusión educativa define la finalidad del sistema; el \DUA\ aporta la lógica pedagógica para anticipar la diversidad; la accesibilidad digital establece las condiciones mínimas de participación en entornos mediados por tecnología; y la educación basada en datos ofrece los recursos analíticos necesarios para transformar información educativa en evidencia interpretable. Esta articulación conduce a un modelo de diagnóstico entendido como proceso continuo, contextualizado y orientado a apoyos pedagógicos.

Con base en esta arquitectura conceptual, se presenta ahora la lógica funcional del sistema, es decir, la manera en que el modelo teórico se traduce en procesos de captura de datos, análisis, generación de diagnósticos, recomendaciones pedagógicas y mejora continua.

\begin{figure}[p]
\centering
\begin{adjustbox}{max totalsize={\textwidth}{0.90\textheight},center}
\begin{tikzpicture}[
font=\small,
node distance=12mm and 12mm,
source/.style={
  draw,
  rounded corners=2mm,
  align=center,
  inner sep=7pt,
  text width=3.8cm
},
process/.style={
  draw,
  rounded corners=2mm,
  align=center,
  inner sep=8pt,
  text width=8.8cm
},
module/.style={
  draw,
  rounded corners=2mm,
  align=center,
  inner sep=7pt,
  text width=4.0cm
},
result/.style={
  draw,
  rounded corners=2mm,
  align=center,
  inner sep=7pt,
  text width=4.2cm
},
arrow/.style={
  -{Stealth[length=2.8mm]},
  line width=0.9pt
},
frame/.style={
  draw,
  dashed,
  rounded corners=2mm,
  inner sep=8pt
}
]

% ======================================================
% FUENTES DE DATOS
% ======================================================
\node[source] (data1) {Datos pedagógicos\\
Actividades\\
Evaluaciones\\
Desempeño observado};

\node[source, right=of data1] (data2) {Datos de interacción\\
Uso de plataforma\\
Trazas digitales\\
Secuencias de navegación};

\node[source, right=of data2] (data3) {Datos contextuales\\
Características del estudiante\\
Condiciones de acceso\\
Entorno educativo};

% ======================================================
% CAPA DE INTEGRACIÓN
% ======================================================
\node[process, below=16mm of data2] (data_layer) {\textbf{Capa de integración de datos educativos}\\
Convergencia de múltiples fuentes de información para construir una base analítica del diagnóstico educativo inclusivo};

\draw[arrow] (data1.south) -- ++(0,-5mm) -| ([xshift=-28mm]data_layer.north);
\draw[arrow] (data2.south) -- (data_layer.north);
\draw[arrow] (data3.south) -- ++(0,-5mm) -| ([xshift=28mm]data_layer.north);

% ======================================================
% ANALÍTICA
% ======================================================
\node[process, below=14mm of data_layer] (analytics) {\textbf{Motor de analítica del aprendizaje}\\
Procesamiento, organización e interpretación de datos para identificar patrones, trayectorias, indicadores y señales relevantes del proceso de aprendizaje};

\draw[arrow] (data_layer) -- (analytics);

% ======================================================
% SALIDAS INTERMEDIAS
% ======================================================
\node[module, below left=14mm and 12mm of analytics] (profile) {Perfil de aprendizaje\\
Ritmos\\
Persistencia\\
Patrones de participación};

\node[module, below=14mm of analytics] (barriers) {Identificación de barreras\\
De acceso\\
De participación\\
De comprensión y desempeño};

\node[module, below right=14mm and 12mm of analytics] (strengths) {Fortalezas y apoyos\\
Recursos disponibles\\
Potencialidades observadas\\
Necesidades de acompañamiento};

\draw[arrow] (analytics.south) -- ++(0,-6mm) -| (profile.north);
\draw[arrow] (analytics.south) -- (barriers.north);
\draw[arrow] (analytics.south) -- ++(0,-6mm) -| (strengths.north);

% ======================================================
% RECOMENDACIONES
% ======================================================
\node[process, below=16mm of barriers] (recommendations) {\textbf{Generación de recomendaciones pedagógicas}\\
Traducción de la evidencia analítica en orientaciones docentes alineadas con los principios del \DUA, la accesibilidad y la atención a la diversidad};

\draw[arrow] (profile.south) -- ++(0,-5mm) -| ([xshift=-24mm]recommendations.north);
\draw[arrow] (barriers.south) -- (recommendations.north);
\draw[arrow] (strengths.south) -- ++(0,-5mm) -| ([xshift=24mm]recommendations.north);

% ======================================================
% PANEL DOCENTE
% ======================================================
\node[process, below=14mm of recommendations] (dashboard) {\textbf{Panel docente de apoyo a decisiones}\\
Visualización integrada de diagnósticos, indicadores, barreras identificadas, fortalezas observadas y recomendaciones pedagógicas interpretables};

\draw[arrow] (recommendations) -- (dashboard);

% ======================================================
% RESULTADOS
% ======================================================
\node[result, below left=16mm and 18mm of dashboard] (out1) {\textbf{Diagnóstico inclusivo}\\
Dinámico\\
Contextualizado\\
Pedagógicamente interpretable};

\node[result, below right=16mm and 18mm of dashboard] (out2) {\textbf{Decisiones pedagógicas}\\
Basadas en evidencia\\
Ajustes didácticos\\
Apoyos diferenciados};

\draw[arrow] (dashboard.south) -- ++(0,-5mm) -| (out1.north);
\draw[arrow] (dashboard.south) -- ++(0,-5mm) -| (out2.north);

% ======================================================
% MEJORA CONTINUA
% ======================================================
\node[module, right=16mm of dashboard] (feedback) {Mejora continua\\
Retroalimentación docente\\
Refinamiento del modelo\\
Iteración DBR};

\draw[arrow] (dashboard.east) -- (feedback.west);
\draw[arrow] (feedback.north) |- ([xshift=10mm]analytics.east);

% ======================================================
% MARCO FUNCIONAL
% ======================================================
\node[frame, fit=(data_layer) (analytics) (profile) (barriers) (strengths) (recommendations) (dashboard)] (mainframe) {};

\node[above=6mm of mainframe, align=center, font=\small\bfseries]
{Arquitectura funcional del sistema de diagnóstico educativo inclusivo};

\end{tikzpicture}
\end{adjustbox}
\caption{Arquitectura funcional del sistema de diagnóstico educativo inclusivo basado en datos.}
\label{fig:arquitectura-funcional-diagnostico}
\end{figure}

La Figura \ref{fig:arquitectura-funcional-diagnostico} representa la traducción operativa del modelo teórico. En ella se observa que el sistema parte de la integración de múltiples fuentes de datos educativos, las procesa mediante analítica del aprendizaje y produce salidas interpretables relacionadas con perfiles, barreras, fortalezas y necesidades de apoyo. Posteriormente, esta información se transforma en recomendaciones pedagógicas alineadas con el \DUA\ y se presenta en un panel docente de apoyo a decisiones. Finalmente, la retroalimentación recogida en la práctica permite refinar el sistema, en coherencia con la lógica iterativa de la Investigación Basada en el Diseño.

En conjunto, ambas figuras permiten comprender que el sistema propuesto no se limita a una arquitectura técnica, sino que articula fundamentos pedagógicos, criterios de accesibilidad, principios éticos y procedimientos analíticos para construir un diagnóstico educativo inclusivo más robusto, contextualizado y útil para la práctica docente.

\printbibliography[heading=subbibliography,title={Referencias (Capítulo II)}]

\end{refsection}